Ko ustvarimo našo ChArUco ploščo, jo bomo natisnili in fotografirali. Posneti moramo okoli 10 slik, ki morajo biti iz različnih kotov. Te slike je treba dati v mapo \texttt{calibdata/} kot jpg. Nato je mogoče zagnati \texttt{calib.py} in generirati bosta \texttt{intrinsic.npy} in \texttt{distCoeff.npy}, ki sta notranja matrika in koeficienti popačenja tudi napaka ponovne projekcije(re-projection error). \cite{calibtut}

\subsection{Parametri plošče in ChArUco detektor}
Na vrhu \texttt{calib.py} najprej deklariramo parametre naše plošče. \texttt{width} in \texttt{height} sta število vrstic in stolpcev, ki smo jih določili prej. \texttt{square\_len} in \texttt{marker\_len} sta realni velikosti kvadratov in oznak, merjeni v metrih. \texttt{dict} je slovarska številka ChArUco vzorca. Z njimi lahko ustvarimo objekt \texttt{aruco\_CharucoBoard}.

\begin{lstlisting}[language=Python, caption={CharucoBoard nastavitve},captionpos=b]
width = 7
height = 10
square_len = 27.2/1000
marker_len = 14/1000
dict = 12
    
aruco_dict = cv.aruco.getPredefinedDictionary(dict)
board_size = (width, height)
board = cv.aruco.CharucoBoard(board_size, square_len, marker_len, aruco_dict)
\end{lstlisting}

\noindent Nato potrebujemo objekt \texttt{aruco\_CharucoDetector}. Za njegovo ustvarjanje potrebujemo board objekt in parametre. Parametre vključujejo CharucoParameters, DetectorParameters in RefineParameters. Če želimo uporabiti izboljšano zaznavo plošče, moramo v CharucoParameters nastaviti \texttt{tryRefineMarkers} na \texttt{true}. S tem bomo zaznali več markerjev.

\begin{lstlisting}[language=Python, caption={CharucoDetector nastavitve},captionpos=b]
refparams = cv.aruco.RefineParameters(minRepDistance = 10.,
									   errorCorrectionRate = 3.,
									   checkAllOrders=True)
    
charucoparms = cv.aruco.CharucoParameters()
charucoparms.tryRefineMarkers = True
    
charuco_detector = cv.aruco.CharucoDetector(board, 
                                            refineParams=refparams,
                                            charucoParams=charucoparms)
\end{lstlisting}


\subsection{Branje slik}
Z uporabo \texttt{glob} dobimo imena datotek slik v mapi \texttt{calibdata/} in za vsako od njih izvedemo zanko. Z \texttt{cv.imread()} preberemo sliko, nato pa jo z \texttt{cv.cvtColor(img, cv.COLOR\_BGR2GRAY)} pretvorimo v sivine. Poenostavitev slike je možna tudi s threshold, vendar povzroči večjo napako ponovne projekcije.

\begin{lstlisting}[language=Python, caption={Branje slik},captionpos=b]
images = glob.glob('calibdata/*.jpg')
    
for fname in images:    
    img = cv.imread(fname)
    gray = cv.cvtColor(img, cv.COLOR_BGR2GRAY)
\end{lstlisting}
\subsection{Zaznavanje vogalov}
\sloppy
S funkcijo \texttt{detectBoard()} objekta detektorja charuco lahko zaznamo vogale in oznake. Položaje vogalov dodamo matriki \texttt{imgpoints} in pripnemo najdene vogale matriki \texttt{objpoints}. Ta niz shranjuje, kje so točke na plošči (npr. ((0,0,0), (1,0,0), (2,0,0))). Najdene vogale lahko tudi narišemo z \texttt{cv.aruco.drawDetectedCornersCharuco(gray, charuco\_corners, charuco\_ids)} in prikažemo rezultat z \texttt{cv.imshow()}.


\begin{lstlisting}[language=Python, caption={Zaznavanje vogalov},captionpos=b]
    charuco_corners, charuco_ids, marker_corners, marker_ids = charuco_detector.detectBoard(gray)
    
    imgpoints.append(charuco_corners)

    objp = np.zeros(((height-1)*(width-1),3), np.float32)
    objp[:,:2] = np.mgrid[0:(width-1),0:(height-1)].T.reshape(-1,2)
    
    objtemp = []
    for id in charuco_ids:
        objtemp.append(objp[int(id)])
  
    objpoints.append(np.array(objtemp))

    cv.aruco.drawDetectedCornersCharuco(gray, charuco_corners, charuco_ids)
    cv.imshow('detected corners', gray)
    
\end{lstlisting}

\subsection{Kalibracija kamere in shranjevanje}
Zdaj lahko izvedemo kalibracijo kamere z \texttt{cv.calibrateCamera(objpoints, imgpoints, (w, h), None, None)}. Ta funkcija vrne notranjo matriko \texttt{mtx}, koeficiente popačenja \texttt{dist}, rotacijske vektorje in translacijske vektorje (zunanji parametri, edinstveni za vsako sliko) \texttt{rvec, tvec}. Notranjo matriko in koeficiente popačenja lahko shranimo za kasnejšo uporabo s funkcijo numpy \texttt{np.save()}. 

\begin{lstlisting}[language=Python, caption={Kalibracija kamere in shranjevanje},captionpos=b]
ret,mtx,dist,rvecs,tvecs = cv.calibrateCamera(objpoints, imgpoints,
                                                    (w, h),
                                                    None, None)

np.save('intrinsic.npy', mtx)
np.save('distCoeff.npy', dist)
\end{lstlisting}

\subsection{Napaka ponovne projekcije}
Z vrednostmi, ki smo jih dobili s kalibracijo kamere, lahko izračunamo napako ponovne projekcije, kar je dobro merilo za oceno, kako dobra je naša kalibracija. Bližje kot je 0, tem bolje.

\begin{lstlisting}[language=Python, caption={Računanje napaka ponovne projekcije},captionpos=b]
mean_error = 0
for i in range(len(objpoints)):
    imgpoints2, _ = cv.projectPoints(objpoints[i], rvecs[i], tvecs[i], mtx, dist)
    error=cv.norm(imgpoints[i],imgpoints2,cv.NORM_L2)/len(imgpoints2)
    mean_error += error

print("Total error: {}".format(mean_error/len(objpoints)))
\end{lstlisting}